\section{The finite-weight ideal on the exceptional rank-three locus}

Fix a coefficient discrete valuation ring $O$, its residue field $k$,
a frame $F$ satisfying $F.\mathrm{Ok}(C_0)$, and its rank-two and
rank-three normalized spaces $S_2,S_3$. In rank $i$ write $A_i$ for
the entry algebra, $B_i\subset A_i$ for the closed coefficient
algebra, $u_i$ for the universal representation, and $\lambda_i$
for the ideal generated by $\mathfrak m_O$ and by
$\psi_{v,a}(g)-1$ for $g\in I_v$ with $\Psi(\operatorname{loc}_v g)=1$.
The condition $\mathrm{FinLam}(I)$ means
$\mathfrak m_{B_i}\subseteq\sqrt{I+\lambda_i}$.
A prime is off the two triangular loci if some entry above and
some entry below the diagonal survive in its quotient. In rank
three, a type-E prime is such a prime containing the four sets
of entries $01,10,12,21$. We transfer a rank-two ideal to every
one of these primes.

% BEGIN C67E1
\paragraph{C67E1. Removing the middle character in the original order.}
Let $E\subset A_3$ be generated by all four sets of entries just
listed. Suppose $E\ne A_3$ and put
\[
      C=B_3/(E\cap B_3),\qquad K=A_3/E,\qquad
      p:B_3\twoheadrightarrow C,\quad q:A_3\twoheadrightarrow K .
\]
By contraction $C\hookrightarrow K$. The ring $C$ is a nonzero
complete Noetherian local $O$-algebra of residue field $k$, by RTC6
and passage to a quotient. Over $K$ the universal matrices preserve the
middle coordinate and the outer coordinates $(0,2)$ separately.
Write $a(g)$ for the middle entry and $\sigma(g)$ for the outer
matrix. Both are multiplicative. The anti-diagonal form and the
marked diagonal equation restrict to those of rank two. We verify
the ordered local equation before using the rank-two universal space.

Write $r_x=\sigma(\operatorname{loc}_v x)$ and $c_i(x)=\psi_{v,i}(x)$
after specialization. The separator condition in $F.\mathrm{Ok}$
gives $s$ with $\Psi(\operatorname{loc}_v s)^2\ne1$.
Set $D=r_s-c_1(s)I$ and $P_x=r_x-c_0(x)I$.
The determinant of $D$ is the image of a unit of $C$. Indeed its
two diagonal factors reduce to $z-1,z^{-1}-1$, where
$z=\Psi(\operatorname{loc}_v s)$, and the product of the opposite
entries $02,20$ belongs to the maximal ideal of the coefficient ring
by the balanced-monomial condition. Similarly both
$a(\operatorname{loc}_v s)-c_0(s)$ and
$a(\operatorname{loc}_v s)-c_2(s)$ come from units of $C$.
The scalar middle block of the three-factor ordinary relation at
$(s,x,s)$ gives $a(\operatorname{loc}_v x)=c_1(x)$.

The outer block at $(x,s,y)$ gives
$P_xD(r_y-c_2(y)I)=0$ for every $x,y$. In particular
$P_xD^2=(c_2(s)-c_1(s))P_xD$. Applying Cayley--Hamilton to $D$ gives
\[
    (\det D)P_x=(\operatorname{tr}D-c_2(s)+c_1(s))P_xD .
\]
Multiplying on the right by $r_y-c_2(y)I$ and cancelling the unit
$\det D$ gives the ordered rank-two product.
Thus $\sigma$ together with the two characters $c_0,c_2$ is
a normalized rank-two representation over the actual quotient $K$.
The calculation works over rings with nilpotents and requires no
splitting of local flags over a field.
% END C67E1

% BEGIN C67E2
\paragraph{C67E2. Compatible maps on entries and on coefficients.}
The representation of C67E1 satisfies the cone conditions of APD9
over $C\hookrightarrow K$. Here are the verifications. Universal
diagonal entries and ordinary characters take their values in $B_3$,
with the prescribed reductions. Applying $p$ gives their values in
$C$. The two outer weights in rank three are $1,-1$, exactly the
rank-two weights. Thus every balanced rank-two entry monomial maps
to the corresponding balanced rank-three monomial; the presence of
an off-diagonal factor still puts it in $\mathfrak m_C$. Each of
the four outer entry spans is generated over $C$ by the images of a
finite generating set for its entry span over $B_3$, by RTC5--RTC9
on the universal space. For the continuity clause let $W$ be the
finite weight piece of an entry and $N\subset W$ its entry span.
For every $r$, Artin--Rees gives $t$ such that
$\mathfrak m_{B_3}^t W\cap N\subseteq\mathfrak m_{B_3}^r N$.
Nearby entry differences belong to the left side by RTC9 and the
definition of $N$. Their images belong to $\mathfrak m_C^r$ times
the new entry span, as required. The outer local characters
descend in the same way and retain the two residual powers of
$\Psi$. These are the diagonal, monomial, finite generation,
continuity and ordinary-character clauses in APD9.

Applying APD10 to this precise cone gives $O$-algebra homomorphisms
\[
  \xi:A_2\longrightarrow K,\qquad \beta:B_2\longrightarrow C,
  \qquad \xi_*u_2=\sigma,\qquad \xi|_{B_2}=\beta
                    \quad\hbox{inside }K .
\]
APD10 realizes the cone by its graded vertex and returns entry values
in $K$ and coefficient values in $C$. In particular the formulas
give $\beta(\psi_{v,i}^{(2)}(g))=p(\psi_{v,2i}^{(3)}(g))$ for
$i=0,1$, including all inertia elements in the definition of $\lambda_i$.
% END C67E2

% BEGIN C67E3
\paragraph{C67E3. The middle character after fixing the weights.}
Let $P$ be a prime of $B_3$ containing $E\cap B_3$ and
$\lambda_3$. The domain $D=B_3/P$ is complete Noetherian local,
of characteristic $\ell$ and residue field $k$. The middle entries
$a(g)$ come from $C$, since each is diagonal. Their images in $D$
define a character $\alpha:G_0\to D^\times$ because the middle
block is multiplicative in $K$ and $C\hookrightarrow K$.
The universal continuity of RTC9 descends through the local quotient.
The middle residual weight is zero, so $\alpha\equiv1\bmod\mathfrak m_D$.
The original polarization $r(g^*)=Jr(g)^tJ$ gives
$\alpha(g^*)=\alpha(g)$. The middle-block identity in C67E1 and
the generators of $\lambda_3$ give $\alpha(\operatorname{loc}_v g)=1$
on $I_v\cap\ker(\Psi\circ\operatorname{loc}_v)$.

For each $N\ge1$, reduce $\alpha$ modulo $\mathfrak m_D^N$. Its image
is a finite abelian $\ell$-group. It is finite by local constancy
and compactness of $G_0$, and its elements have $\ell$-power order
by $(1+y)^{\ell^n}=1+y^{\ell^n}$ for $y\in\mathfrak m_D$.
Its kernel is open. This character is fixed by the star operation
and kills all the marked inertia generators $s'_a$ of ACC9 by the
preceding paragraph. The second detector in ACC9 consequently
bounds its exponent by the same $H>0$ for every $N$.
Krull intersection gives $\alpha(g)^H=1$ in $D$. Write
$H=\ell^bH_0$ with $\ell\nmid H_0$. Since $D$ is a domain,
$(\alpha(g)^{H_0}-1)^{\ell^b}=0$ implies $\alpha(g)^{H_0}=1$.
The geometric sum dividing $X^{H_0}-1$ by $X-1$ evaluates at
$\alpha(g)$ to a unit with residue $H_0$. Hence $\alpha(g)=1$
for every $g$. This proves the required triviality only after passage
to $D$, retaining the middle character on the original exceptional
quotient.
% END C67E3

% BEGIN C67E4
\paragraph{C67E4. Finiteness of the transferred weight ideal.}
Let $I_2\subseteq B_2$ satisfy $\mathrm{FinLam}(I_2)$, and put
\[
                I_3=p^{-1}\!\left(C\cdot\beta(I_2)\right).
\]
For any prime $P\supset I_3+\lambda_3$ we prove
$\mathfrak m_{B_3}\subseteq P$. In fact $\ker p\subset I_3$.
Thus $\bar P=p(P)$ is prime in $C$ and $p^{-1}\bar P=P$.
The prime $\beta^{-1}\bar P$ contains $I_2$ and contains $\lambda_2$
by the coefficient and outer-character formulas in C67E2.
Rank-two finiteness forces $\mathfrak m_{B_2}\subseteq\beta^{-1}\bar P$.
The image of $B_2$ in $D=B_3/P$ therefore consists of scalars from
$k$, since the structural map $O\to B_2/\mathfrak m_{B_2}=k$
is surjective. The image of $O$ in $D$ identifies canonically with
this same $k$ because $\lambda_3$ kills $\mathfrak m_O$.

Every weight-zero product of rank-three entries is zero over $K$
if it uses an odd entry. Otherwise its factors are middle entries
and entries in the outer block. Removing the middle factors leaves
a weight-zero monomial in the rank-two entries, since the middle
weight difference is zero and the outer weights agree. This latter
monomial belongs to $B_2$ by definition; its value in $K$ is its
image under $\beta$ in $C$. The middle factors belong to $C$ and
become one in $D$ by C67E3. Injectivity $C\hookrightarrow K$ hence
identifies the image of every balanced monomial in $D$ with a
scalar. The outer ordinary characters are scalars by C67E2; the
middle ones are one by C67E1 and C67E3.

By RTC6 and RTC8 the polynomial algebra on these monomials and
ordinary characters is dense in $B_3$ for its own maximal-ideal
topology: the defining closure in the complete deformation ring
induces the same topology on its degree-zero part. Its image in
$D$ is dense because the quotient map is continuous and
surjective. That image is contained in the finite subfield $k$,
a closed subset of the Hausdorff ring $D$. Thus $D=k$, proving
$\mathfrak m_{B_3}\subseteq P$. Intersecting over all such primes
gives $\mathrm{FinLam}(I_3)$.
% END C67E4

% BEGIN C67E5
\paragraph{C67E5. A single ideal covering every type-E prime.}
Assume also the original rank-two support hypothesis:
$I_2\subset Q\cap B_2$ for every type-I prime $Q$ of $A_2$.
Take any type-E prime $P'$ of $A_3$. As $E\subseteq P'$,
$P'/E$ is a prime of $K$. Its inverse image $Q$ under $\xi$
is type I in rank two. Indeed $P'$ omits at least one strictly
lower and one strictly upper entry. The odd entries all belong
to $P'$, so an entry $20$ and an entry $02$ survive. These entries
are exactly the lower and upper off-diagonal entries of the rank-two
corner. Rank two has no extra exceptional exclusion in its type-I
predicate. Therefore $I_2\subset Q\cap B_2$. The compatibility
$\xi|_{B_2}=\beta$ yields $C\beta(I_2)\subset (P'/E)\cap C$
and hence $I_3\subset P'\cap B_3$. The same $I_3$ works for every
$P'$. If $E=A_3$ there are no type-E primes and $I_3=B_3$ satisfies
$\mathrm{FinLam}$ directly. This proves the original $\mathrm{l53}$
for all $F.\mathrm{Ok}(C_0)$ and arbitrary $S_2,S_3$.

In particular, once a rank-three ideal $M$ satisfies
$\mathrm{FinLam}(M)$ and covers the type-I primes, ACP18 applies
to $M$ and this $I_3$. Their intersection satisfies
$\mathrm{FinLam}$ and covers every off-triangular prime, giving
the hypothesis $\mathrm{OffFin}$ for ACP11. The construction of
$I_3$ itself uses the separate rank-two ideal and its support
hypothesis.
% END C67E5
